%% 08_conclusion.tex

\section{Conclusion}
\label{sec:conclusion}

We presented the first paired-data comparison of narrative RAG, naive structured FHIR
RAG, and resource-aware structured FHIR RAG on specialty-medication adherence-indicator
question answering, where information content is held constant across representations
by constructing all three retrieval substrates from the same FHIR bundles.
Narrative RAG achieves the highest exact-match accuracy (40.6\% vs.\ 33.4--35.3\%)
under the current 512-token output budget, a result substantially attributable to
the reasoning-truncation artefact that penalises structured systems with denser
contexts.  Temporal-anchor failures are a structural weakness of the narrative format
that persists independent of budget.  In the feature-extraction arm, structured FHIR
features substantially outperform both narrative-derived and retrieval-trace features
for adherence classification, demonstrating the representation dissociation that is
the paper's central empirical claim: winning natural-language QA and winning downstream
adherence-outcome ML are different objectives, and the same representation does not
simultaneously optimise both.

The dataset (200 patients, 13,800 questions, five PSP adherence-indicator families),
code (three retrieval systems, evaluation harness, feature extractors), and analysis
scripts are released as a reusable benchmark for the specialty-pharmacy QA research
community.

Future directions include: (1) sweeping ANSWER\_MAX\_TOKENS to isolate the architecture
effect from the budget artefact; (2) evaluating on real de-identified PSP data post-REB;
(3) testing robustness across multiple answer LLMs and embedding models; (4) extending
the question bank to additional specialty-medication classes and real-world edge cases
(infusion regimens, variable-interval dosing, bridge therapies); and (5) constructing
richer synthetic adherence labels to stress-test the feature-extraction arm beyond
the perfectly-scheduled synthetic scenario.
